\documentclass[10pt,a4paper]{amsart}
\usepackage{amsmath,amssymb,amsthm,graphicx}
\usepackage{verbatim}
\newenvironment{reference}{\comment}{\endcomment}
%\newenvironment{reference}{}{}
\newenvironment{slogan}{\comment}{\endcomment}
\newenvironment{history}{\comment}{\endcomment}

% For commutative diagrams we use Xy-pic
\usepackage[all]{xy}

% We use 2cell for 2-commutative diagrams.
\xyoption{2cell}
% Theorem environments.
%
\theoremstyle{plain}
\newtheorem{theorem}[subsection]{Theorem}
\newtheorem{proposition}[subsection]{Proposition}
\newtheorem{lemma}[subsection]{Lemma}

\theoremstyle{definition}
\newtheorem{definition}[subsection]{Definition}
\newtheorem{example}[subsection]{Example}
\newtheorem{exercise}[subsection]{Exercise}
\newtheorem{situation}[subsection]{Situation}

\theoremstyle{remark}
\newtheorem{remark}[subsection]{Remark}
\newtheorem{remarks}[subsection]{Remarks}

\numberwithin{equation}{subsection}

% Macros
%
\def\lim{\mathop{\mathrm{lim}}\nolimits}
\def\colim{\mathop{\mathrm{colim}}\nolimits}
\def\Spec{\mathop{\mathrm{Spec}}}
\def\Hom{\mathop{\mathrm{Hom}}\nolimits}
\def\Ext{\mathop{\mathrm{Ext}}\nolimits}
\def\SheafHom{\mathop{\mathcal{H}\!\mathit{om}}\nolimits}
\def\SheafExt{\mathop{\mathcal{E}\!\mathit{xt}}\nolimits}
\def\Sch{\mathit{Sch}}
\def\Mor{\mathop{\mathrm{Mor}}\nolimits}
\def\Ob{\mathop{\mathrm{Ob}}\nolimits}
\def\Sh{\mathop{\mathit{Sh}}\nolimits}
\def\NL{\mathop{N\!L}\nolimits}
\def\CH{\mathop{\mathrm{CH}}\nolimits}
\def\proetale{{pro\text{-}\acute{e}tale}}
\def\etale{{\acute{e}tale}}
\def\QCoh{\mathit{QCoh}}
\def\Ker{\mathop{\mathrm{Ker}}}
\def\Im{\mathop{\mathrm{Im}}}
\def\Coker{\mathop{\mathrm{Coker}}}
\def\Coim{\mathop{\mathrm{Coim}}}

% Boxtimes
%
\DeclareMathSymbol{\boxtimes}{\mathbin}{AMSa}{"02}

%
% Macros for moduli stacks/spaces
%
\def\QCohstack{\mathcal{QC}\!\mathit{oh}}
\def\Cohstack{\mathcal{C}\!\mathit{oh}}
\def\Spacesstack{\mathcal{S}\!\mathit{paces}}
\def\Quotfunctor{\mathrm{Quot}}
\def\Hilbfunctor{\mathrm{Hilb}}
\def\Curvesstack{\mathcal{C}\!\mathit{urves}}
\def\Polarizedstack{\mathcal{P}\!\mathit{olarized}}
\def\Complexesstack{\mathcal{C}\!\mathit{omplexes}}
% \Pic is the operator that assigns to X its picard group, usage \Pic(X)
% \Picardstack_{X/B} denotes the Picard stack of X over B
% \Picardfunctor_{X/B} denotes the Picard functor of X over B
\def\Pic{\mathop{\mathrm{Pic}}\nolimits}
\def\Picardstack{\mathcal{P}\!\mathit{ic}}
\def\Picardfunctor{\mathrm{Pic}}
\def\Deformationcategory{\mathcal{D}\!\mathit{ef}}
\usepackage{cleveref}
\crefname{originallemma}{original lemma}{original lemmas}
\Crefname{originallemma}{Original lemma}{Original lemmas}
\makeatletter
\newlabel{lemma-induced-and-quotient-strict}{{0124}{source}{}{sourcepointer-1}{}}
\newlabel{lemma-induced-and-quotient-strict@cref}{{[originallemma][0][]0124}{[1][0][]source}{}{}{}}
\makeatother
\title{The composition of strict morphisms of filtered objects is strict under an injectivity or surjectivity hypothesis}
\author{The Stacks Project Authors}
\date{stacks-project:homology.tex:local-snapshot-20261012; GNU FDL 1.2 or later}
\begin{document}
\maketitle
\noindent\textit{Editorial note.} All delivered mathematical text is the author's own, copied from the pinned Stacks Project chapter file homology.tex; no statement or proof step was rewritten, reordered or removed.\par\medskip
\noindent\textit{Editorial note.} The proof is the author's complete argument for all three assertions, including the counterexample; each displayed computation is followed by the justification of its individual steps, and the closing sentence is a conclusion rather than a placeholder.\par\medskip
\noindent\textit{Editorial note.} Every labelled reference occurring inside the counted spans is carried in assembly.external\_references as a non-load-bearing navigation pointer (same-chapter ones at their definition site, cross-chapter ones at the invocation inside the counted span), so no load-bearing mathematics is replaced by a stub. No bibliography citation occurs inside the counted spans.\par\medskip
\noindent\textit{Editorial note.} The same-name check for this card scanned every staged Stacks chapter of this repository -- all sources/stacks-*/pinned-original-source.tex files, the card's own chapter excluded -- and the labels found are recorded in dedup.same\_name\_cross\_chapter on the card.\par\medskip
\noindent\textit{Editorial note.} The preamble is the chapter's own native preamble excerpt already used for this source: the comment-type reference, slogan and history environments (native preamble lines 8-13), the Xy-pic setup (native preamble lines 14-18) and the native theorem declarations with the mathematical macros (native preamble lines 154-217, the actual last line of the file).\par\medskip
\noindent\textit{Allowed standard prerequisites (editorial).}\begin{itemize}
\item Strict morphisms of filtered objects. Same chapter, the definition of a strict morphism of filtered objects. Supplies the condition on the induced filtration of the image of a morphism that all three assertions are about.
\item Induced and quotient filtrations are strict. Same chapter, Lemma lemma-induced-and-quotient-strict. Supplies that the inclusion and the quotient used in the counterexample are strict morphisms of filtered objects, so that the counterexample really lies inside the hypotheses of part (1).
\end{itemize}\medskip
\ifdefined\hypertarget\hypertarget{sourcepointer-1}{}\fi
\noindent\textit{Source reference.} Original lemma 0124: sources/stacks-homology/pinned-original-source.tex, source lines 5020--5020. Navigation pointer invoked at proof line 5050; the counted span contains the pointer only and the cited result is not reproduced inside it, so the pointer is registered as a non-load-bearing navigation stub in the established form for Stacks material. No bibliography citation is involved.\par\medskip
\begin{lemma}
\label{lemma-composition-strict}
Let $\mathcal{A}$ be an abelian category. Let $f : A \to B$, $g : B \to C$
be strict morphisms of filtered objects.
\begin{enumerate}
\item In general the composition $g \circ f$ is not strict.
\item If $g$ is injective, then $g \circ f$ is strict.
\item If $f$ is surjective, then $g \circ f$ is strict.
\end{enumerate}
\end{lemma}

\begin{proof}
Let $B$ a vector space over a field $k$ with basis $e_1, e_2$, with the
filtration $F^nB = B$ for $n < 0$, with $F^0B = ke_1$, and $F^nB = 0$ for
$n > 0$. Now take $A = k(e_1 + e_2)$ and $C = B/ke_2$ with filtrations
induced by $B$, i.e., such that $A \to B$ and $B \to C$ are strict
(Lemma \ref{lemma-induced-and-quotient-strict}).
Then $F^n(A) = A$ for $n < 0$ and $F^n(A) = 0$ for $n \geq 0$. 
Also $F^n(C) = C$ for $n \leq 0$ and $F^n(C) = 0$ for $n > 0$.
So the (nonzero) composition $A \to C$ is not strict.

\medskip\noindent
Assume $g$ is injective. Then
\begin{align*}
g(f(F^pA)) & = g(f(A) \cap F^pB) \\
& = g(f(A)) \cap g(F^p(B)) \\
& = (g \circ f)(A) \cap (g(B) \cap F^pC) \\
& = (g \circ f)(A) \cap F^pC.
\end{align*}
The first equality as $f$ is strict, the second because $g$ is injective,
the third because $g$ is strict, and the fourth because
$(g \circ f)(A) \subset g(B)$.

\medskip\noindent
Assume $f$ is surjective. Then
\begin{align*}
(g \circ f)^{-1}(F^iC) & = f^{-1}(F^iB + \Ker(g)) \\
& = f^{-1}(F^iB) + f^{-1}(\Ker(g)) \\
& = F^iA + \Ker(f) + \Ker(g \circ f) \\
& = F^iA + \Ker(g \circ f)
\end{align*}
The first equality because $g$ is strict, the second because $f$ is
surjective, the third because $f$ is strict, and the last because
$\Ker(f) \subset \Ker(g \circ f)$.
\end{proof}

\end{document}
